\documentclass[10pt]{article}

\usepackage{hyperref}
\usepackage{xcolor}
\usepackage{calc}
\usepackage{graphicx}
\usepackage{tikz}
\usepackage{fontspec}
\usepackage{fontawesome5}
\usepackage{titlesec}
\usepackage{enumitem}
\usepackage{fancybox}
\usepackage{xeCJK}

\hypersetup{hidelinks}
\urlstyle{same}

%%%%%%%%%%%%%%%%%%%%
% 设置
%%%%%%%%%%%%%%%%%%%%

\setlength{\parindent}{0pt}
\setlength{\parskip}{0pt}
\pagenumbering{gobble}
\setlist[itemize]{
    leftmargin=1.35em,
    itemsep=0.12em,
    topsep=0.18em,
    parsep=0pt,
    partopsep=0pt
}
\setlist[enumerate]{leftmargin=*}
\renewcommand{\arraystretch}{1.12}
\linespread{1.08}

\titleformat{\section}
  {\Large\bfseries\raggedright}
  {}{0em}
  {}
  [{\color{secondarycolor}\titlerule}]
\titlespacing*{\section}{0pt}{0.55em}{0.28em}

\usepackage[
    a4paper,
    left=1.1cm,
    right=1.1cm,
    top=0.75cm,
    bottom=0.85cm,
    nohead
]{geometry}

% 字体设置
\setCJKmainfont[
    Path=fonts/,
    Extension=.otf,
    BoldFont=*-Bold,
]{NotoSerifSC}

% 自定义颜色
\definecolor{primarycolor}{RGB}{200, 60, 60}
\definecolor{secondarycolor}{RGB}{200, 68, 28}

\newlength{\iconwidth}
\setlength{\iconwidth}{1.5em}

\newcommand{\projectheading}[3]{%
    \noindent\makebox[\textwidth][s]{%
        \makebox[0pt][l]{{\large \textbf{#1}}}%
        \hfill
        \makebox[0pt][c]{{\normalsize \textbf{#2}}}%
        \hfill
        \makebox[0pt][r]{#3}%
    }%
}

%%%%%%%%%%%%%%%%%%%%
% 文章内容
%%%%%%%%%%%%%%%%%%%%

% 联系方式
\newcommand{\contact}{
    \footnotesize
    \textcolor{white}{
        \faEnvelope \quad 邮箱：\href{mailto:xxx@qq.com}{250xxxxxxx@qq.com}
        \hspace{3.2em}
        \faPhone \quad 电话：198xxxxxxxx
        \hspace{3.2em}
        \faGithub \quad 主页：\href{https://github.com/Dithob}{github.com/Dithob}
    }
}

\begin{document}

    % 顶部联系人栏
    \begin{tikzpicture}[remember picture, overlay]
        \node[anchor=north, inner sep=0pt](headercontacts) at (current page.north){
            \includegraphics[width=\paperwidth]{images/foot.png}
        };
        \node[anchor=center] at(headercontacts.center){\contact};
    \end{tikzpicture}
    \vspace{-1.15em}

    %%%%%%%%%%%%%%%%%%%%
    % 个人信息与教育背景
    %%%%%%%%%%%%%%%%%%%%

    \begin{minipage}[t]{0.79\textwidth}
        \section[个人信息]{\makebox[\iconwidth][c]{\color{primarycolor}{\faAddressCard}}\quad 个人信息}
        {\small
        \begin{tabular*}{\textwidth}{@{\extracolsep{\fill}}lll}
            \textbf{姓\qquad 名}：无绝 & \textbf{出生年月}：2002.10 & \textbf{籍\qquad 贯}：湖北咸宁 \\
            \textbf{学\qquad 校}：重庆邮电大学 & \textbf{英语水平}：CET-6 & \textbf{求职方向}：前端开发
        \end{tabular*}
        }
        % \vspace{0.75em}

        \section[教育背景]{\makebox[\iconwidth][c]{\color{primarycolor}{\faGraduationCap}}\quad 教育背景}
        {\large \textbf{重庆邮电大学}}，计算机技术，\textit{硕士} \hfill 2024.09--2027.06
        \begin{itemize}
            \item \textbf{荣誉奖项}：特等学业奖学金 $\times$ 1、二等学业奖学金 $\times$ 1
        \end{itemize}

        \vspace{0.1em}
        {\large \textbf{重庆邮电大学}}，计算机科学与技术，\textit{本科} \hfill 2020.09--2024.07
        \begin{itemize}
            \item \textbf{荣誉奖项}：三等学业奖学金；GPA: 3.28/4.0（专业前 25\%）
        \end{itemize}

    \end{minipage}
    \hfill
    \begin{minipage}[t]{0.14\textwidth}
        \vspace{-0.35em}
        \setlength{\fboxsep}{0pt}
        \doublebox{\includegraphics[width=\linewidth]{images/xjpic.jpg}}
    \end{minipage}

    % \vspace{0.45em}

    %%%%%%%%%%%%%%%%%%%%
    % 专业技能
    %%%%%%%%%%%%%%%%%%%%

    \section[专业技能]{\makebox[\iconwidth][c]{\color{primarycolor}{\faWrench}}\quad 专业技能}
    \begin{itemize}
        \item \textbf{前端基础与框架}：熟悉 HTML5、CSS3、JavaScript、TypeScript，掌握 Vue3、Vuex/Pinia 等技术栈，能够完成页面组件化开发、状态管理、路由组织与复杂交互实现。
        \item \textbf{组件与可视化}：熟悉 Ant Design Vue、Element 等组件库，能够封装表单、列表、弹窗、搜索筛选等通用业务组件；熟悉 ECharts，可实现基础数据看板与可视化展示。
        \item \textbf{工程化与联调}：掌握 Git、Docker、Postman、JMeter、FoxAPI 等工具，熟悉前后端接口联调、请求/响应数据排查、环境部署与基础性能测试流程。
        \item \textbf{后端与数据基础}：了解 Django、SpringBoot、MyBatis 等后端框架，掌握 MySQL、SQLite、Redis 的基本使用，能够参与接口设计、数据库建模和前后端数据流打通。
        \item \textbf{AI 辅助开发}：熟悉 GitHub Copilot、Claude Code、Codex、Qoder 等 AI 编程工具，能够结合 AI 对话、SSE 流式响应和实时预览能力开发 AI 应用型前端项目。
    \end{itemize}

    % \vspace{0.45em}

    %%%%%%%%%%%%%%%%%%%%
    % 项目经历
    %%%%%%%%%%%%%%%%%%%%

    \section[项目经历]{\makebox[\iconwidth][c]{\color{primarycolor}{\faProjectDiagram}}\quad 项目经历}

    \projectheading{云端智能图像协同系统}{前端开发}{2025.06--2025.08}
    \begin{itemize}
        \item \textbf{项目目标}：面向公共图库、个人私有图库与企业团队共享场景，开发支持图片检索、多维分析、管理员审核与 AI 实时协同编辑的图像管理平台。
        \item \textbf{权限与任务状态}：基于 Vue Router 为路由节点配置 \texttt{meta.access} 字段，并结合 \texttt{beforeEach} 全局守卫实现页面级访问控制；对接后端 AI 图像处理模型，针对“AI 扩图”等高耗时任务采用“任务派发--定时轮询”机制追踪异步执行状态。
        \item \textbf{协同编辑与可视化}：基于 WebSocket 事件驱动架构开发多人协同编辑能力，引入编辑锁与冲突检测机制，降低并发编辑下的状态覆写风险；使用 ECharts 封装分组条形图、词云图等数据组件，并借助 Vue3 响应式能力动态更新图表配置。
    \end{itemize}

    \vspace{0.25em}

    \projectheading{智能应用生成平台}{前端开发}{2025.09--2025.12}
    \begin{itemize}
        \item \textbf{项目目标}：面向非技术用户的 AI 应用生成平台，支持通过多轮 AI 对话生成网站，并提供实时响应式预览与可视化修改能力，降低轻量 Web 应用搭建门槛。
        \item \textbf{流式交互与预览通信}：采用原生 \texttt{EventSource} 消费后端 SSE 流式接口，实现 AI 对话打字机输出效果，并动态展示模型工具调用轨迹；使用 \texttt{postMessage} 打通主应用与 \texttt{iframe} 预览站点之间的跨域通信，支撑实时预览与编辑状态同步。
        \item \textbf{可视化编辑与组件复用}：通过向子 \texttt{iframe} 动态注入 DOM 监听脚本，实现“点击即修改”的可视化编辑能力；基于 Ant Design Vue 封装应用卡片、信息弹窗等业务组件，通过 \texttt{props} 和 \texttt{emits} 进行父子组件通信，提高页面开发复用率。
    \end{itemize}

    \vspace{0.25em}

    \projectheading{校园二手交易平台}{项目负责人}{2025.03--2025.06}
    \begin{itemize}
        \item \textbf{项目目标}：面向校园二手交易场景，开发支持商品发布、商品浏览、关键词搜索与交易流程的信息平台，覆盖用户、商品与交易等核心实体管理。
        \item \textbf{页面与交互实现}：负责商品列表、商品详情、商品发布和搜索结果等核心页面开发，围绕列表渲染、条件筛选、表单提交和状态提示组织前端交互逻辑；基于响应式布局适配不同屏幕尺寸。
        \item \textbf{模板复用与数据流打通}：使用 Django 模板引擎抽取页面公共结构，实现导航、商品卡片、列表区域等模块复用；参与后端 API 设计与 SQLite 实体关系建模，打通商品数据展示、提交、查询和交易状态流转。
    \end{itemize}

    % \vspace{0.45em}

    %%%%%%%%%%%%%%%%%%%%
    % 实习经历
    %%%%%%%%%%%%%%%%%%%%

    \section[实习经历]{\makebox[\iconwidth][c]{\color{primarycolor}{\faBriefcase}}\quad 实习经历}

    \projectheading{重庆啄木鸟网络科技有限公司}{算法实习生（AI 应用研发）}{2025.06--2025.09}
    \begin{itemize}
        \item \textbf{核心职责}：参与 AI 应用研发，负责模型调研、工作流搭建、接口设计与测试、后端功能开发及前后端联调，参与智能产品识别、小啄 GPT 智能回复、Agent 文档解析与联网搜索等模块。
        \item \textbf{智能家电识别与自动入库系统}：使用 Docker/Miniconda 统一部署 OCR/视觉大模型，接入 vLLM 优化推理服务；基于 Dify 搭建“识别--联网补全--字段校验--向量去重--入库”工作流，并通过接口封装为前端提供结构化产品信息、补全结果和去重结果数据。
    \end{itemize}

    % \vspace{0.45em}

    %%%%%%%%%%%%%%%%%%%%
    % 奖项证书
    %%%%%%%%%%%%%%%%%%%%

    \section[奖项证书]{\makebox[\iconwidth][c]{\color{primarycolor}{\faTrophy}}\quad 奖项证书}
    \begin{itemize}
        % \item \textbf{证书}：CET-6
        \item \textbf{竞赛}：2025 “互联网+”大学生创新创业大赛银奖；2025 首届重庆市 AI 大模型创新应用大赛省级三等奖；2023 第十四届蓝桥杯 C/C++ 组省赛三等奖；2023 重庆市高校数据库应用竞赛校级二等奖
    \end{itemize}

\end{document}
